\section{ĐỀ SỐ \thesection}%ĐỀ 1 
\setcounter{bt}{0}
%%Câu 1.
\begin{bt} Thực hiện các phép tính:
	\begin{listEX}[2]
		\item $-2-(-12)+(-18)$;
		\item $126^{15}: 126^{14}-26 \cdot(40-38)^{2}$.
	\end{listEX}
	\loigiai{a) $-8$ \hfill b) $22$.
	}
\end{bt}

%%Câu 2.
\begin{bt} Nhiệt độ ban đầu trong một kho lạnh là $-20^{\circ} \mathrm{C}$. Do yêu cầu bảo quản hàng hóa, người quản lý kho tiếp tục giảm $3^{\circ} \mathrm{C}$. Sau đó, do hàng hóa được lấy bớt nên quản lý kho tăng nhiệt độ thêm $5^{\circ} \mathrm{C}$. Hỏi sau khi tăng, nhiệt độ hiện tại của kho lạnh là bao nhiêu?
	
	\loigiai{$-18^\circ$C.}
\end{bt}


%%Câu 3.
\begin{bt}  Một đội nghi thức của một trường Trung học cơ sở A có khoảng $100$ đến $150$ đội viên đã tham gia hội thi "Nghi thức Đội". Khi tham gia biểu diễn, đội xếp đội hình $10$ hàng, $12$ hàng, $15$ hàng đều vừa đủ. Hỏi đội nghi thức của trường Trung học cơ sở A có bao nhiêu đội viên tham gia hội thi?
	\loigiai{$120$ đội viên.}
\end{bt}


%%Câu 4.
\begin{bt} \textbf{ }\\
	\begin{minipage}[h]{0.5\linewidth}
		Biểu đồ tranh ở hình bên cho ta thông tin về loại nhạc cụ yêu thích của các bạn học sinh khối lớp 6 của trường Trung học cơ sở A. Biết rằng không có học sinh không thích nhạc cụ.
		\begin{listEX}
			\item Số lượng học sinh yêu thích nhạc cụ Organ nhiều hơn số học sinh yêu thích nhạc cụ trống là bao nhiêu?
			\item Tính tổng số học sinh của trường Trung học cơ sở A.
		\end{listEX}
	\end{minipage}
	\begin{minipage}{0.5\linewidth}
		\begin{center}
			\def\icon{\huge \quarterNote}
			\def\hicon{\huge \eighthNote}
			\def\vetranh#1{\foreach \dem in{1,...,#1}{\icon \ }}{
				\begin{tabular}{|c|p{4cm}<{}|}\hline 
					\bf Loại nhạc cụ&\multicolumn{1}{c|}{\bf Số học sinh yêu thích}\\ \hline 
					Organ &\vetranh{4}\ \hicon \\ \hline 
					Ghi ta &\vetranh{5}\ \hicon  \\ \hline 
					Kèn &\vetranh{2}\ \hicon \\ \hline 
					Trống &\vetranh{4} \\ \hline 
					Sáo &\vetranh{2}\ \hicon \\ \hline 
				\end{tabular}\\
				(\icon \normalsize = 10 học sinh;        \hicon \normalsize = 5 học sinh)}	
		\end{center}
		
	\end{minipage}
	\loigiai{a) $5$ học sinh \hfill b) $190$ học sinh.
	}
\end{bt}


%%Câu 5.
\begin{bt} 
	\immini[thm]{
		Một bồn hoa hình chữ nhật (như hình vẽ), có \\
		$AB = HF = 4$ m, $BC = EG = 3$ m. Người ta trồng hoa hồng bên trong phần đất hình thoi $HEFG$, phần đất còn lại trồng hoa mười giờ.
		\begin{listEX}
			\item Tính diện tích đất trồng hoa hồng?
			\item Tính diện tích đất trồng hoa mười giờ?
		\end{listEX}
	}{
		\begin{tikzpicture}[scale=1,line width=0.8pt,line join=round,>=stealth]
			\def\a{4} \def\b{3} 
			\path (0,0) coordinate (D)
			(0,\b) coordinate (A)
			(\a,\b) coordinate (B)
			(\a,0) coordinate (C)
			($(A)!0.5!(B)$) coordinate (E)
			($(B)!0.5!(C)$) coordinate (F)
			($(C)!0.5!(D)$) coordinate (G)
			($(A)!0.5!(D)$) coordinate (H)
			;
			\draw (A)--(B)--(C)--(D)--(A)
			(E)--(F)--(G)--(H)--(E)
			;
			\draw[<->] (A)++(90:0.6)--++(0:\a)node[midway,above]{4 m};
			\draw[<->] (B)++(0:0.6)--++(-90:\b)node[midway,right]{3 m};
			\foreach \x/ \goc in {A/135,B/45,C/-45,D/-135,E/90,F/0,G/-90,H/180}  \fill (\x) circle (0pt)        
			($(\x)+(\goc:3.5mm)$) node {$\x$};
		\end{tikzpicture}
	}
	\loigiai{a) $6$ m$^2$ \hfill b) $6$ m$^2$.}
\end{bt}


%%Câu 6. 
\begin{bt} Chứng minh: $\left(7+7^{2}+7^{3}+7^{4}+\ldots+7^{48}+7^{49}+7^{50}+7^{51}\right)$ chia hết cho $399$.
	
	\loigiai{HD: Nhận thấy $(7+7^{2}+7^{3})$ $\vdots$ $399$\\
		Dùng phương pháp kết hợp và tính chất phân phối của phép nhân đối với phép cộng.}
\end{bt}
